% ID: SEL-D0019-OB-E03
% ESTADO: validada. S020, cotejada con examen y criterios oficiales.
\textbf{a)} Al sumar $P(A)+P(B)$ se cuenta dos veces la intersección.
Por tanto $P(B)=P(A\cup B)-P(A)+P(A\cap B)=\boxed{0.6}$.

\textbf{b)} La misma identidad da $P(C\cap D)=0.4+0.3-0.5=0.2$.
Así $P(C\cap D^c)=P(C)-P(C\cap D)=0.2$. Entre los casos en que no
ocurre $D$, la proporción correspondiente es
$P(C\mid D^c)=0.2/(1-0.3)=\boxed{2/7}$.

\textbf{c)} Los complementarios de sucesos independientes también son
independientes. Así $P(E^c\cap F^c)=(1-0.6)(1-0.8)=\boxed{0.08}$.
